\section{Introduction}
\label{sec:intro}

An information system is not merely a body of code; it is a system of work in which organisational
rules are encoded so that routine decisions can be made consistently. When a lender approves a loan, a retailer prices an order, or an insurer settles a claim, the software enacts a business rule that someone, somewhere, wrote down. The distance between the rule as intended and the rule as enacted is
an old concern of the Information Systems (IS) field, whose literature has long treated the alignment
between business strategy and information technology as a continuing managerial problem rather than a
one-off design decision \citep{henderson1993strategic}, and which distinguishes the quality of the
information an artefact carries from the quality of the system that carries it
\citep{delone2003model}. This paper is about a specific, measurable, and surprisingly neglected way in
which that distance grows.

Behaviour-Driven Development (BDD) was proposed precisely to close it. A requirement is written as a
concrete scenario in a structured natural language (Gherkin); each step of the scenario is bound to
code; and the scenario is executed as a test. The promise is that the specification \emph{is} the
test, so the specification cannot silently diverge from the implementation. BDD is now widely used in
industry, and empirical work has confirmed its adoption in open-source projects: in a study of
50{,}000 popular repositories, approximately 27\,\% used a BDD framework, rising to 68\,\% among Ruby
projects \citep{zampetti2020demystifying}.

The promise has a structural weakness. A Gherkin scenario is only an oracle if its steps actually
\emph{check} something that could fail. Scenarios can be tagged \texttt{@wip} and skipped; assertions
inside step definitions can be removed, loosened, or written so that the observed value is compared
against a value computed by the very production code under test. When that happens the scenario still
runs, still passes, still contributes to coverage --- and no longer detects the change it exists to
detect. The test has become a mirror of the defect rather than a check on it. Coverage metrics cannot
see this, because coverage measures reach, not discrimination.

The consequence for practice is uncomfortable: an organisation may believe it holds an executable
specification of its business rules while in fact holding a set of scripts that would pass under
almost any implementation. The consequence for IS research is an open measurement question. Is there
an instrument that tells an organisation, with acceptable effort, whether its specification layer
still functions as an oracle --- and whether that ability is decaying?

\subsection{The gap this paper addresses}

The literature tells us a great deal about BDD's adoption \citep{zampetti2020demystifying}, its
scenario quality \citep{chandorkar2022gherkin,binamungu2020characterising}, the maintainability burden
of keeping scenarios and code in step \citep{binamungu2018maintaining,oliveira2019howtoevaluate}, and
its benefits as reported by practitioners \citep{couto2023characterization}. It tells us that mutation
analysis is the standard instrument for measuring a suite's discriminating power, and that mutation
scores track real fault detection in a way coverage does not \citep{Just2014AreMutants}. It also tells
us that tests decay: disabled tests are common and often never re-enabled \citep{Kim2021Disabled},
test cases have a measurable half-life \citep{Feldt2014GrowOld}, and tests can pass while checking
nothing \citep{Delplanque2019RottenGreen,zhi2013inadequateassertions}.

What the literature does not provide is the intersection. A recent systematic mapping of 166 BDD
papers published between 2006 and 2021 found no treatment of mutation analysis, and explicitly named
an ``acute shortage of metrics for measuring various aspects of BDD specifications''
\citep{binamungu2023mapping}. The closest existing work applied mutation analysis to a BDD acceptance
suite exactly once, in a single industrial robotic-system proof of concept with hand-selected domain
mutants \citep{santos2024aat4irs}. To our knowledge, no study has measured, across a corpus of public
Cucumber-style suites, how much discriminating power the specification layer contributes
\emph{relative to the rest of the test suite}, nor how that contribution changes over repository
history. This paper is an attempt to build the instrument for that measurement, and an honest report
of what happened when we pointed it at the public ecosystem.

\subsection{Research questions}

The study is organised around three questions that separate \emph{whether we can measure} from
\emph{what we find}:

\begin{description}
\item[RQ1 (Measurement validity).] Can specification drift be operationalised as a reproducible
      metric set, and are the resulting measurements reliable and tool-feasible across ecosystems?
\item[RQ2 (Oracle strength).] What share of covered, domain-layer mutants does the BDD suite
      \emph{alone} kill, and how much does it add beyond the repository's other tests?
\item[RQ3 (Association and erosion).] Do pre-specified static drift indicators associate with lower
      BDD-only kill rates, and does drift accumulate over repository history?
\end{description}

A word on vocabulary, because it matters for honesty. We use \emph{oracle strength} for the share of
covered domain-layer mutants killed by the specification layer. This is a \emph{snapshot} property.
\emph{Drift} is a \emph{change over time}: a suite whose oracle strength is low at one commit may
simply always have been weak. We are careful throughout not to infer drift from a single measurement,
and the term \emph{specification drift} is used here as a definition we introduce, not as an
established construct in the literature (Section~\ref{sec:related} explains the established
relatives: architectural drift and erosion \citep{PerryWolf1992,Rosik2010ArchDrift}, code decay
\citep{Izurieta2012Decay}).

\subsection{Contributions}

\begin{enumerate}
\item \textbf{An instrument.} A five-step protocol (\texttt{docs/PROTOCOL.md}) that a team can apply
      to its own suite: delimit the domain layer, measure oracle strength in BDD-only and full-suite
      modes with mandatory coverage conditioning, compute seven drift indicators, test for
      accumulation over history, and interpret under explicit rules that forbid the most common
      misreadings (Section~\ref{sec:design}, Section~\ref{sec:method}).
\item \textbf{Tool-feasibility evidence, established empirically rather than assumed.} Mutation
      analysis driven by the BDD runner itself, with per-mutant coverage and per-scenario
      attribution, is demonstrated in both the Java ecosystem (PIT with the Cucumber-JVM
      JUnit-Platform engine) and the JavaScript/TypeScript ecosystem (Stryker with the official
      Cucumber runner). We report the exact pinned versions and the observed mutation counts.
\item \textbf{A pre-specified, quantified account of the population.} From a frame of \Nframe{}
      repositories enumerated without a popularity-ranked search, we report eligibility, build and
      execution attrition with categorised reasons, and show that the binding constraint on this
      literature is infrastructural rather than statistical.
\item \textbf{A negative result and its implications.} The claim that public BDD suites can serve as
      the empirical basis for a corpus-scale study of oracle quality does not survive contact with
      the ecosystem. We quantify why, and we state what a future study would have to do instead.
\end{enumerate}

\subsection{What this paper does not claim}

We do not claim that BDD is ineffective, nor that public BDD suites are of poor quality. We do not
establish that BDD causes better or worse oracles: the design is observational and observational
designs cannot support that claim. We do not measure defects or field failures; mutation score is a
\emph{proxy} for the ability to detect a change in business-rule behaviour, not a count of real bugs.
And we do not report inter-rater reliability for the two judgement-coded indicators, because no
independent human rater was available; we say so explicitly rather than reporting a number we could
not honestly produce (Section~\ref{sec:threats}).

\subsection{Organisation}

Section~\ref{sec:background} states the problem in IS terms. Section~\ref{sec:related} positions the
work against prior literature and states precisely what is and is not new.
Section~\ref{sec:design} presents the design-science framing and the artefact.
Section~\ref{sec:method} gives the pre-specified method: frame, eligibility, metrics, mutation
protocol and statistics. Section~\ref{sec:results} reports results, including the measured attrition.
Section~\ref{sec:discussion} discusses implications for research and practice.
Section~\ref{sec:threats} argues each threat to validity and quantifies it where possible.
Section~\ref{sec:conclusion} concludes.

